\documentclass[a4paper,12pt]{article}
\usepackage{amsmath}
\usepackage{geometry}
\usepackage{fancyhdr}
\usepackage{multicol}

\geometry{top=2cm, bottom=2cm, left=2.5cm, right=2.5cm}
\pagestyle{fancy}
\fancyhf{}
\rhead{Esc. 4-117 Ejército de los Andes}
\lhead{Soluciones Evaluación – Termodinámica}
\cfoot{\thepage}

\begin{document}

\section*{Resolución – Evaluación 1}

\begin{enumerate}
\item \textbf{Ley de Charles:}
\[
V_1 = 1{,}8\,L,\quad T_1 = 280\,K,\quad T_2 = 350\,K
\]
\[
V_2 = V_1 \cdot \frac{T_2}{T_1} = 1{,}8 \cdot \frac{350}{280} = 2{,}25\,L
\]

\item \textbf{Ley de Gay-Lussac:}
\[
P_1 = 1{,}0\,atm,\quad T_1 = 10^\circ C = 283\,K,\quad T_2 = 373\,K
\]
\[
P_2 = P_1 \cdot \frac{T_2}{T_1} = 1{,}0 \cdot \frac{373}{283} \approx 1{,}32\,atm
\]

\item \textbf{Ecuación general:}
\[
P_1 = 2{,}0\,atm,\quad V_1 = 3{,}0\,L,\quad T_1 = 300\,K\\
P_2 = 1{,}2\,atm,\quad T_2 = 330\,K
\]
\[
V_2 = \frac{P_1 V_1 T_2}{P_2 T_1}
= \frac{2{,}0 \cdot 3{,}0 \cdot 330}{1{,}2 \cdot 300} = \frac{1980}{360} = 5{,}5\,L
\]

\item \textbf{Gas ideal (cálculo de volumen):}
\[
n = 0{,}5\,mol,\quad T = 273\,K,\quad P = 1\,atm,\quad R = 0{,}082
\]
\[
V = \frac{nRT}{P} = \frac{0{,}5 \cdot 0{,}082 \cdot 273}{1} \approx 11{,}2\,L
\]

\item \textbf{Ley de Dalton (presiones parciales):}
\[
n_{tot} = 0{,}4 + 0{,}3 + 0{,}3 = 1{,}0\,mol,\quad P_{total} = 1{,}2\,atm
\]
\[
P_{O_2} = 0{,}4 \cdot 1{,}2 = 0{,}48\,atm,\quad P_{N_2} = 0{,}36\,atm,\quad P_{H_2} = 0{,}36\,atm
\]

\item \textbf{Ley de Amagat (volúmenes parciales):}
\[
V_{total} = 2{,}0 + 1{,}5 + 2{,}5 = 6{,}0\,L
\]
\[
\%V_{O} = 33{,}3\%,\quad \%V_{N} = 25\%,\quad \%V_{H} = 41{,}7\%
\]

\end{enumerate}

\newpage

\section*{Resolución – Evaluación 2}

\begin{enumerate}
\item \textbf{Ley de Charles:}
\[
V_1 = 2{,}1\,L,\quad T_1 = 290\,K,\quad T_2 = 360\,K
\]
\[
V_2 = V_1 \cdot \frac{T_2}{T_1} = 2{,}1 \cdot \frac{360}{290} \approx 2{,}60\,L
\]

\item \textbf{Ley de Gay-Lussac:}
\[
P_1 = 0{,}9\,atm,\quad T_1 = 288\,K,\quad T_2 = 368\,K
\]
\[
P_2 = 0{,}9 \cdot \frac{368}{288} \approx 1{,}15\,atm
\]

\item \textbf{Ecuación general:}
\[
P_1 = 1{,}3\,atm,\quad V_1 = 5{,}0\,L,\quad T_1 = 280\,K\\
P_2 = 1{,}8\,atm,\quad T_2 = 330\,K
\]
\[
V_2 = \frac{1{,}3 \cdot 5{,}0 \cdot 330}{1{,}8 \cdot 280} = \frac{2145}{504} \approx 4{,}25\,L
\]

\item \textbf{Gas ideal (cálculo de moles):}
\[
P = 2{,}0\,atm,\quad V = 6{,}0\,L,\quad T = 300\,K
\]
\[
n = \frac{PV}{RT} = \frac{12{,}0}{0{,}082 \cdot 300} \approx 0{,}49\,mol
\]

\item \textbf{Ley de Dalton (presiones parciales):}
\[
n_{tot} = 2 + 1 + 1 = 4\,mol,\quad P_{total} = 2{,}4\,atm
\]
\[
P_{CO_2} = 1{,}2\,atm,\quad P_{CH_4} = 0{,}6\,atm,\quad P_{O_2} = 0{,}6\,atm
\]

\item \textbf{Ley de Amagat (volúmenes parciales):}
\[
V_{total} = 3{,}0 + 1{,}0 + 2{,}0 = 6{,}0\,L
\]
\[
\%V_A = 50\%,\quad \%V_B = 16{,}7\%,\quad \%V_C = 33{,}3\%
\]

\end{enumerate}

\end{document}
